\begin{register}{H}{HP\_SYSTEM\_EXTERNAL\_DEVICE\_ENCRYPT\_DECRYPT\_CONTROL\_REG}{0x0000}\label{regdesc:HPSYSTEMEXTERNALDEVICEENCRYPTDECRYPTCONTROLREG}
\regfield{\rotatebox{15}{(reserved)}}{28}{4}{0000000000000000000000000000}%
\regfield{\rotatebox{15}{HP\_SYSTEM\_ENABLE\_DOWNLOAD\_MANUAL\_ENCRYPT}}{1}{3}{{0}}%
\regfield{\rotatebox{15}{HP\_SYSTEM\_ENABLE\_DOWNLOAD\_G0CB\_DECRYPT}}{1}{2}{{0}}%
\regfield{\rotatebox{15}{reserved}}{1}{1}{{0}}%
\regfield{\rotatebox{15}{HP\_SYSTEM\_ENABLE\_DOWNLOAD\_DB\_ENCRYPT}}{1}{1}{{0}}%
\regfield{\rotatebox{15}{HP\_SYSTEM\_ENABLE\_SPI\_MANUAL\_ENCRYPT}}{1}{0}{{0}}%
\reglabel{Reset}\regnewline%
\begin{regdesc}\begin{reglist}
\label{fielddesc:HPSYSTEMENABLESPIMANUALENCRYPT}\item [HP\_SYSTEM\_ENABLE\_SPI\_MANUAL\_ENCRYPT] 配置是否在 SPI Boot 模式下开启 MSPI XTS 手动加密。\\
0：关闭\\
1：开启\\
(R/W)
\label{fielddesc:HPSYSTEMENABLEDOWNLOADDBENCRYPT}\item [HP\_SYSTEM\_ENABLE\_DOWNLOAD\_DB\_ENCRYPT] 配置是否在 Joint Download Boot 模式下启用自动加密。\\
0：禁用\\
1：启用\\
(R/W)
\label{fielddesc:HPSYSTEMENABLEDOWNLOADG0CBDECRYPT}\item [HP\_SYSTEM\_ENABLE\_DOWNLOAD\_G0CB\_DECRYPT] 配置是否在下载 Boot 模式下开启 MSPI XTS 自动解密。\\
0：关闭\\
1：开启\\
(R/W)
\label{fielddesc:HPSYSTEMENABLEDOWNLOADMANUALENCRYPT}\item [HP\_SYSTEM\_ENABLE\_DOWNLOAD\_MANUAL\_ENCRYPT] 配置是否在下载 Boot 模式下开启 MSPI XTS 手动加密。\\
0：关闭\\
1：开启\\
(R/W)
\end{reglist}\end{regdesc}
\end{register}


\begin{register}{H}{HP\_SYSTEM\_SEC\_DPA\_CONF\_REG}{0x0008}\label{regdesc:HPSYSTEMSECDPACONFREG}
\regfield{(reserved)}{29}{3}{00000000000000000000000000000}%
\regfield{HP\_SYSTEM\_SEC\_DPA\_CFG\_SEL}{1}{2}{{0}}%
\regfield{HP\_SYSTEM\_SEC\_DPA\_LEVEL}{2}{0}{{0x0}}%
\reglabel{Reset}\regnewline%
\begin{regdesc}\begin{reglist}
\label{fielddesc:HPSYSTEMSECDPALEVEL}\item [HP\_SYSTEM\_SEC\_DPA\_LEVEL] 配置是否开启防 DPA 攻击。仅当 \hyperref[fielddesc:HPSYSTEMSECDPACFGSEL]{HP\_SYSTEM\_SEC\_DPA\_CFG\_SEL} 置 1 时可用。\\
0：关闭\\
1-3：开启。值越大，安全等级越高，防 DPA 攻击的能力就越强。不过，硬件加密加速器的计算开销也随之增加。\\
(R/W)
\label{fielddesc:HPSYSTEMSECDPACFGSEL}\item [HP\_SYSTEM\_SEC\_DPA\_CFG\_SEL] 配置选择 \hyperref[fielddesc:HPSYSTEMSECDPALEVEL]{HP\_SYSTEM\_SEC\_DPA\_LEVEL} 或是 \hyperref[fielddesc:EFUSESECDPALEVEL]{EFUSE\_SEC\_DPA\_LEVEL} (eFuse 寄存器) 来控制 DPA 等级。\\
0：选择 \hyperref[fielddesc:EFUSESECDPALEVEL]{EFUSE\_SEC\_DPA\_LEVEL} \\
1：选择 \hyperref[fielddesc:HPSYSTEMSECDPALEVEL]{HP\_SYSTEM\_SEC\_DPA\_LEVEL}\\
(R/W)
\end{reglist}\end{regdesc}
\end{register}


\begin{register}{H}{HP\_SYSTEM\_CPU\_PERI\_TIMEOUT\_CONF\_REG}{0x000C}\label{regdesc:HPSYSTEMCPUPERITIMEOUTCONFREG}
\regfield{(reserved)}{14}{18}{00000000000000}%
\regfield{HP\_SYSTEM\_CPU\_PERI\_TIMEOUT\_PROTECT\_EN}{1}{17}{{1}}%
\regfield{HP\_SYSTEM\_CPU\_PERI\_TIMEOUT\_INT\_CLEAR}{1}{16}{{0}}%
\regfield{HP\_SYSTEM\_CPU\_PERI\_TIMEOUT\_THRES}{16}{0}{{0xffff}}%
\reglabel{Reset}\regnewline%
\begin{regdesc}\begin{reglist}
\label{fielddesc:HPSYSTEMCPUPERITIMEOUTTHRES}\item [HP\_SYSTEM\_CPU\_PERI\_TIMEOUT\_THRES] 配置总线访问 CPU 外设寄存器的超时阈值，与时钟域的时钟周期数相对应。 (R/W)
\label{fielddesc:HPSYSTEMCPUPERITIMEOUTINTCLEAR}\item [HP\_SYSTEM\_CPU\_PERI\_TIMEOUT\_INT\_CLEAR] 写 1 清除超时中断。(WT)
\label{fielddesc:HPSYSTEMCPUPERITIMEOUTPROTECTEN}\item [HP\_SYSTEM\_CPU\_PERI\_TIMEOUT\_PROTECT\_EN] 配置是否启用访问 CPU 外设寄存器的超时保护。\\
0：关闭\\
1：开启\\
(R/W)
\end{reglist}\end{regdesc}
\end{register}


\begin{register}{H}{HP\_SYSTEM\_CPU\_PERI\_TIMEOUT\_ADDR\_REG}{0x0010}\label{regdesc:HPSYSTEMCPUPERITIMEOUTADDRREG}
\regfield{HP\_SYSTEM\_CPU\_PERI\_TIMEOUT\_ADDR}{32}{0}{{0x000000}}%
\reglabel{Reset}\regnewline%
\begin{regdesc}\begin{reglist}
\label{fielddesc:HPSYSTEMCPUPERITIMEOUTADDR}\item [HP\_SYSTEM\_CPU\_PERI\_TIMEOUT\_ADDR] 表示异常访问的地址信息。(RO)
\end{reglist}\end{regdesc}
\end{register}


\begin{register}{H}{HP\_SYSTEM\_CPU\_PERI\_TIMEOUT\_UID\_REG}{0x0014}\label{regdesc:HPSYSTEMCPUPERITIMEOUTUIDREG}
\regfield{(reserved)}{25}{7}{0000000000000000000000000}%
\regfield{HP\_SYSTEM\_CPU\_PERI\_TIMEOUT\_UID}{7}{0}{{0x0}}%
\reglabel{Reset}\regnewline%
\begin{regdesc}\begin{reglist}
\label{fielddesc:HPSYSTEMCPUPERITIMEOUTUID}\item [HP\_SYSTEM\_CPU\_PERI\_TIMEOUT\_UID] 表示超时发生时的 master ID [4:0] 和 master 权限 [6:5]。详见章节 \ref{mod:permctrl} \textit{\nameref{mod:permctrl}}。该寄存器将在清除超时中断后清除。 (WTC)
\end{reglist}\end{regdesc}
\end{register}


\begin{register}{H}{HP\_SYSTEM\_HP\_PERI\_TIMEOUT\_CONF\_REG}{0x0018}\label{regdesc:HPSYSTEMHPPERITIMEOUTCONFREG}
\regfield{(reserved)}{14}{18}{00000000000000}%
\regfield{HP\_SYSTEM\_HP\_PERI\_TIMEOUT\_PROTECT\_EN}{1}{17}{{1}}%
\regfield{HP\_SYSTEM\_HP\_PERI\_TIMEOUT\_INT\_CLEAR}{1}{16}{{0}}%
\regfield{HP\_SYSTEM\_HP\_PERI\_TIMEOUT\_THRES}{16}{0}{{0xffff}}%
\reglabel{Reset}\regnewline%
\begin{regdesc}\begin{reglist}
\label{fielddesc:HPSYSTEMHPPERITIMEOUTTHRES}\item [HP\_SYSTEM\_HP\_PERI\_TIMEOUT\_THRES] 配置总线访问 HP 外设寄存器的超时阈值，单位为时钟域的时钟周期数。(R/W)
\label{fielddesc:HPSYSTEMHPPERITIMEOUTINTCLEAR}\item [HP\_SYSTEM\_HP\_PERI\_TIMEOUT\_INT\_CLEAR] 写 1 清除超时中断。(WT)
\label{fielddesc:HPSYSTEMHPPERITIMEOUTPROTECTEN}\item [HP\_SYSTEM\_HP\_PERI\_TIMEOUT\_PROTECT\_EN] 配置是否启用访问 HP 外设寄存器的超时保护。\\
0：关闭\\
1：开启\\
(R/W)
\end{reglist}\end{regdesc}
\end{register}


\begin{register}{H}{HP\_SYSTEM\_HP\_PERI\_TIMEOUT\_ADDR\_REG}{0x001C}\label{regdesc:HPSYSTEMHPPERITIMEOUTADDRREG}
\regfield{HP\_SYSTEM\_HP\_PERI\_TIMEOUT\_ADDR}{32}{0}{{0x000000}}%
\reglabel{Reset}\regnewline%
\begin{regdesc}\begin{reglist}
\label{fielddesc:HPSYSTEMHPPERITIMEOUTADDR}\item [HP\_SYSTEM\_HP\_PERI\_TIMEOUT\_ADDR] 表示异常访问的地址信息。(RO)
\end{reglist}\end{regdesc}
\end{register}


\begin{register}{H}{HP\_SYSTEM\_HP\_PERI\_TIMEOUT\_UID\_REG}{0x0020}\label{regdesc:HPSYSTEMHPPERITIMEOUTUIDREG}
\regfield{(reserved)}{25}{7}{0000000000000000000000000}%
\regfield{HP\_SYSTEM\_HP\_PERI\_TIMEOUT\_UID}{7}{0}{{0x0}}%
\reglabel{Reset}\regnewline%
\begin{regdesc}\begin{reglist}
\label{fielddesc:HPSYSTEMHPPERITIMEOUTUID}\item [HP\_SYSTEM\_HP\_PERI\_TIMEOUT\_UID] 表示超时发生时的 master  ID [4:0] 和 master 权限 [6:5]。详见章节 \ref{mod:permctrl} \textit{\nameref{mod:permctrl}}。该寄存器将在清除超时中断后清除。 (WTC)
\end{reglist}\end{regdesc}
\end{register}

\begin{register}{H}{LP\_PERI\_BUS\_TIMEOUT\_CONF\_REG}{0x0010}\label{regdesc:LPPERIBUSTIMEOUTCONFREG}
\regfield{(reserved)}{14}{18}{00000000000000}%
\regfield{LP\_PERI\_BUS\_TIMEOUT\_PROTECT\_EN}{1}{17}{{1}}%
\regfield{LP\_PERI\_BUS\_TIMEOUT\_INT\_CLEAR}{1}{16}{{0}}%
\regfield{LP\_PERI\_BUS\_TIMEOUT\_THRES}{16}{0}{{0xffff}}%
\reglabel{Reset}\regnewline%
\begin{regdesc}\begin{reglist}
\label{fielddesc:LPPERIBUSTIMEOUTTHRES}\item [LP\_PERI\_BUS\_TIMEOUT\_THRES] 配置总线访问 LP 外设寄存器的超时阈值，单位为时钟域的时钟周期数。(R/W)
\label{fielddesc:LPPERIBUSTIMEOUTINTCLEAR}\item [LP\_PERI\_BUS\_TIMEOUT\_INT\_CLEAR] 写 1 清除超时中断。(WT)
\label{fielddesc:LPPERIBUSTIMEOUTPROTECTEN}\item [LP\_PERI\_BUS\_TIMEOUT\_PROTECT\_EN] 配置是否启用访问 LP 外设寄存器的超时保护。\\
0：关闭\\
1：开启\\
(R/W)
\end{reglist}\end{regdesc}
\end{register}


\begin{register}{H}{LP\_PERI\_BUS\_TIMEOUT\_ADDR\_REG}{0x0014}\label{regdesc:LPPERIBUSTIMEOUTADDRREG}
\regfield{LP\_PERI\_BUS\_TIMEOUT\_ADDR}{32}{0}{{0x000000}}%
\reglabel{Reset}\regnewline%
\begin{regdesc}\begin{reglist}
\label{fielddesc:LPPERIBUSTIMEOUTADDR}\item [LP\_PERI\_BUS\_TIMEOUT\_ADDR] 表示异常访问的地址信息。(RO)
\end{reglist}\end{regdesc}
\end{register}


\begin{register}{H}{LP\_PERI\_BUS\_TIMEOUT\_UID\_REG}{0x0018}\label{regdesc:LPPERIBUSTIMEOUTUIDREG}
\regfield{(reserved)}{25}{7}{0000000000000000000000000}%
\regfield{LP\_PERI\_BUS\_TIMEOUT\_UID}{7}{0}{{0x0}}%
\reglabel{Reset}\regnewline%
\begin{regdesc}\begin{reglist}
\label{fielddesc:LPPERIBUSTIMEOUTUID}\item [LP\_PERI\_BUS\_TIMEOUT\_UID] 表示超时发生时的 master  ID [4:0] 和 master 权限 [6:5]。详见章节 \ref{mod:permctrl} \textit{\nameref{mod:permctrl}}。该寄存器将在清除超时中断后清除。 (WTC)
\end{reglist}\end{regdesc}
\end{register}


\begin{register}{H}{HP\_SYSTEM\_DATE\_REG}{0x03FC}\label{regdesc:HPSYSTEMDATEREG}
\regfield{(reserved)}{4}{28}{0000}%
\regfield{HP\_SYSTEM\_DATE}{28}{0}{{0x2412270}}%
\reglabel{Reset}\regnewline%
\begin{regdesc}\begin{reglist}
\label{fielddesc:HPSYSTEMDATE}\item [HP\_SYSTEM\_DATE] 版本控制寄存器。 (R/W)
\end{reglist}\end{regdesc}
\end{register}
